% answers_short.tex —— 2015 年·湖南卷·文科 · 简短答案速查
% 由 tools/make_short_answers.py 自 answers.tex 的 \gksolhead/\gkanswer 提取，
% 并与 2015/2015_数学真题_湖南文_答案要点.md 逐题比对后生成，勿手改（改源 MD 后重跑）。

\begin{gknotice}
\textbf{使用说明}：本节只给最终答案便于快速核对，完整推导见解析版。答案为编写者独立推导，\textbf{非官方标准答案}。
\end{gknotice}

\gkpart{选择题}
\begin{center}\begin{tabular}{ccccccccccc}
\toprule
题号 & 1 & 2 & 3 & 4 & 5 & 6 & 7 & 8 & 9 & 10 \\
\midrule
答案 & D & B & C & A & B & D & C & A & B & A \\
\bottomrule
\end{tabular}\end{center}

\gkpart{填空题}
\begin{itemize}[leftmargin=2.6em]
  \item[\textbf{11.}] \(\{1,2,3\}\)
  \item[\textbf{12.}] \(x^2+(y-1)^2=1\)（即 \(x^2+y^2-2y=0\)）
  \item[\textbf{13.}] \(2\)
  \item[\textbf{14.}] \((0,2)\)
  \item[\textbf{15.}] \(\dfrac{\pi}{2}\)
\end{itemize}

\gkpart{解答题}
\begin{itemize}[leftmargin=2.6em]
  \item[\textbf{16.}] (1) 所有可能结果为 \(A_1a_1\)，\(A_1a_2\)，\(A_1b_1\)，\(A_1b_2\)，\(A_2a_1\)，\(A_2a_2\)，\(A_2b_1\)，\(A_2b_2\)，\(Ba_1\)，\(Ba_2\)，\(Bb_1\)，\(Bb_2\)，共 \(12\) 种；(2) 不正确（中奖概率为 \(\dfrac13\)，小于不中奖概率 \(\dfrac23\)）
  \item[\textbf{17.}] (1) 成立（由 \(a=b\tan A\) 与正弦定理得 \(\sin B=\cos A\)）；(2) \(A=C=\dfrac{\pi}{6}\)，\(B=\dfrac{2\pi}{3}\)
  \item[\textbf{18.}] (1) 成立（\(AE\perp\) 平面 \(B_1BCC_1\)，故平面 \(AEF\perp\) 平面 \(B_1BCC_1\)）；(2) \(\dfrac{\sqrt6}{12}\)
  \item[\textbf{19.}] (1) 成立（由递推式相减得 \(a_{n+2}=3a_n\)）；(2) 当 \(n\) 为偶数时 \(S_n=\dfrac{3(3^{n/2}-1)}{2}\)；当 \(n\) 为奇数时 \(S_n=\dfrac{5\cdot3^{(n-1)/2}-3}{2}\)
  \item[\textbf{20.}] (1) \(\dfrac{y^2}{9}+\dfrac{x^2}{8}=1\)；(2) \(\pm\dfrac{\sqrt6}{4}\)
  \item[\textbf{21.}] (1) 成立（公比 \(q=-\mathrm{e}^{\pi}\)）；(2) \(a\geqslant\dfrac{\sqrt2\pi}{4}\mathrm{e}^{-\pi/4}\)
\end{itemize}
